%! AP Statistics — Unit 4, Lesson 4.7 — Student Packet Cover Sheet
\documentclass[10pt]{article}
\usepackage{apstats-article}
\usepackage{apstats-boxes}
\usepackage{ltablex}
\keepXColumns

\begin{document}

% Full-bleed navy banner
\begin{tikzpicture}[remember picture, overlay]
  \fill[navy] (current page.north west)
    rectangle ([yshift=-0.9in]current page.north east);
\end{tikzpicture}
\vspace{-0.8in}

\begin{center}
  {\color{white}\textbf{\LARGE AP Statistics}}\\[4pt]
  {\color{skymid}\large\textbf{Unit 4: Probability, Random Variables, and Probability Distributions}}\\[6pt]
  {\color{white}\textbf{\large Lesson 4.7 \quad Introduction to Random Variables and Probability Distributions}}
\end{center}

\vspace{0.22in}
\namedateperiod
\vspace{0.18in}

\begin{learningtargetbox}
\small
\begin{enumerate}[leftmargin=1.5em, itemsep=5pt]
  \item \ldots{} represent a probability distribution for a discrete random variable as a
        table, histogram, or formula and verify that it is valid.
  \item \ldots{} build a cumulative probability distribution and use it to answer
        probability questions.
  \item \ldots{} interpret a probability distribution by describing its shape, center,
        and spread in context.
\end{enumerate}
\end{learningtargetbox}

\vspace{0.14in}

\begin{tocbox}
\small
\renewcommand{\arraystretch}{1.5}
\begin{tabularx}{\linewidth}{@{} c l X r @{}}
\rowcolor{sky}
\textbf{\#} & \textbf{Component} & \textbf{Description} & \textbf{Score} \\
\midrule
1 & Warm-Up        & Spiral review: basic probability rules              & \blank{1.2cm} \\
2 & Guided Notes   & Random variables; distribution tables \& histograms & \blank{1.2cm} \\
3 & Group Activity & Building and interpreting distributions (3 tiers)   & \blank{1.2cm} \\
4 & Exit Ticket    & Siblings distribution: validity and interpretation  & \blank{1.2cm} \\
5 & Homework       & Distribution tables, cumulative probabilities       & \blank{1.2cm} \\
\midrule
  & \multicolumn{2}{r}{\textbf{Total}} & \blank{1.2cm} \\
\end{tabularx}
\end{tocbox}

\vspace{0.14in}

\begin{remindbox}
\small
The \textbf{most important idea in Lesson 4.7}: a probability distribution tells the complete
story of a random variable's long-run behavior.

\vspace{6pt}
\begin{tabularx}{\linewidth}{@{} >{\centering\arraybackslash\columncolor{goldbg}}p{0.06\linewidth}
                                    X @{}}
\textbf{\large!} &
\textbf{Two validity conditions --- memorize them.}\enspace
Every probability must be between 0 and 1, and all probabilities must sum to exactly 1. \\[6pt]
\textbf{\large!} &
\textbf{Cumulative means ``at most.''}\enspace
$P(X \le 3)$ means the probability that the random variable is 3 \emph{or less}. \\[6pt]
\textbf{\large!} &
\textbf{Interpretation requires context.}\enspace
Never just restate numbers; say what the shape/center/spread means for the
real-world situation. \\
\end{tabularx}
\end{remindbox}

\end{document}
